\documentclass[10pt,a4paper]{amsart}
\usepackage[margin=19mm]{geometry}
\usepackage{amsmath,amssymb,mathtools}
\usepackage[scr=rsfs,cal=euler]{mathalfa}
\usepackage{xfrac}
\usepackage[unicode,hidelinks,bookmarks=false]{hyperref}
\newcommand{\ie}{i.e.\ }
\newcommand{\cf}{cf.\ }
\newcommand{\resp}{respectively\ }
\newcommand{\Subsection}{\subsection}
\providecommand{\nobreakdash}{\nobreak\hbox{-}}
\setlength{\emergencystretch}{2em}
\multlinegap0pt
\usepackage{amssymb}
\usepackage{dsfont}
\usepackage{float}
\newtheorem{theorem}{Theorem}
\newtheorem{lemma}{Lemma}
\newcommand{\FD}{D}
\newcommand{\N}{\mathbb{N}}
\newcommand{\R}{\mathbb{R}}
\newcommand{\rstar}{r_*}
\newcommand{\xstar}{x_*}
\title{\Large\bf Superconvergence of Centered Finite Difference Approximations}
\author{Mario J. Bencomo, Joseph Igot, Emma Maltes}
\date{Source: arXiv:2609.15831v1 (CC BY 4.0)}
\begin{document}
\maketitle

\noindent\emph{Source and licence.} \Large\bf Superconvergence of Centered Finite Difference Approximations. Version arXiv:2609.15831v1 (CC BY 4.0), distributed by arXiv under the Creative Commons Attribution 4.0 International licence, http://creativecommons.org/licenses/by/4.0/, as recorded in the arXiv licence metadata for this exact version and shown on the abstract page https://arxiv.org/abs/2609.15831v1. Full licence notice: \texttt{licenses/arxiv-2609.15831-CC-BY-4.0.txt}.\par\smallskip

\noindent\textit{Editorial scope.} Counted unit: the article's theorem thm:SCodd (source0.tex lines 436--450). The article's complete original proof of it is delivered with it (source0.tex lines 452--601, one \texttt{proof} environment of 150 lines). No mathematics is written, repaired or re-derived: the statement and the proof are the article's own lines, in the article's own notation, and no retained line is substituted or re-flowed.  Retained uncounted as the source-local context the proof needs: lines 275--288; lines 320--322; lines 1087--1102; lines 1192--1198.  Nothing was omitted from any retained span, so the package carries no omission record.  No retained line contains a citation, so the wrapper carries no bibliography and no bibliography entry.  No retained line contains a figure environment, an image-inclusion command or drawing code, so no image is delivered and the package declares no asset dependency.  Editorial formatting only, in addition: an AMS \texttt{amsart} preamble that restates the article's own declarations and macros used by the retained lines, namely the environments \texttt{theorem}, \texttt{lemma} on separate counters, together with the standard packages those lines need.  The retained environments print in this document as Theorem 1, Lemma 1 in order of appearance, whereas the article prints the same items with the numbering of its own full document.  The complete native source of the article as distributed in the arXiv e-print for this exact version is bundled unchanged as \texttt{sources/210455/source0.tex}.
\begin{theorem}\label{thm:stand_conv}
\text{\rm (Standard Convergence)}\\
Let $f:(a,b)\to\R$ be a sufficiently differentiable function over an open interval $(a,b)$, and $\xstar\in(a,b)$.
Suppose $k,N\in\N$ with $N-k\ge 1$.
Then, for a given fundamental stencil $\{a_n\}_{n=1}^N$ and stencil size $h>0$, there exists a unique set of coefficients $\{c_n\}_{n=1}^N$ such that
\begin{equation}\label{eq:FDappx}
	\FD^k f(\xstar) := \sum_{n=1}^N c_n f(\xstar+a_nh),
\end{equation}
is an FD approximation of $f^{(k)}(\xstar)$ with accuracy order of at least $N-k$.
In other words,
\[
	\left| \FD^k f(\xstar) - f^{(k)}(\xstar) \right| =\mathcal{O}(h^{N-k}).
\]
\end{theorem}

\begin{equation}\label{eq:tilde c}
	c_n = \tilde c_n \frac{k!}{h^k},
\end{equation}

\begin{lemma}\label{lem:BigTheta2}
Suppose the error of an FD approximation is of the form
\begin{equation}\label{eq:err_rs}
	E(h) = |K_1h^r + K_2(h)h^{r+s}|
\end{equation}
where $r,s>0$, and constant $K_1\neq 0$ is independent of $h$.
Assume there exists $0< K_2^*<\infty$ independent of $h$ such that
\[
	|K_2(h)| < K_2^*
\]
for $h>0$ small enough.
It follows that 
\[
	E(h) = \Theta(h^r).
\]
\end{lemma}

\begin{lemma}\label{lem:genVander}
Let $V\in\R^{N\times N}$ be a generalized Vandermonde matrix with entries
\[
	[V]_{mn} = x_n^{z_m}, \quad m=1,2,...,M, \; n=1,2,...,N.
\] 
If $0<x_1 < x_2 < \cdots < x_N$, then $V$ is invertible.
\end{lemma}

\begin{theorem} \label{thm:SCodd}
\text{\rm (Odd-derivative, even-centered stencil case)}\\
Let $f:(a,b)\to\R$ be a sufficiently differentiable function over an open interval $(a,b)$, and $\xstar\in (a,b)$.
Suppose $k=2\ell+1$ and $N=2M$ for some $\ell\in\N_0$ and $M\in\N$, assuming $N-k\ge 1$.
Then, for a given $N$-point, even-centered stencil, with fundamental stencil $\{a_n\}$ and stencil size $h>0$, there exists a unique set of coefficients $\{c_n\}$ such that 
\[
	\FD^k f(\xstar) := \sum_{n=1}^M\Big[ c_{-n}f(\xstar+a_{-n}h)+c_nf(\xstar+a_nh) \Big],
\]
is an FD approximation of $f^{(k)}(\xstar)$ with accuracy order exactly $N-k+1$, assuming $f^{(N+1)}(\xstar)\neq 0$.
In other words, 
\[
	\left| \FD^k f(\xstar) - f^{(k)}(\xstar) \right| = \Theta(h^{N-k+1}).
\]
Furthermore, $\{c_n\}$ must be skew-symmetric.
\end{theorem}

\begin{proof}  
To prove superconvergence of even-centered FD approximations for odd-derivatives we make use of a ``fake'' odd-centered stencil.
Given an $N$-point, even-centered stencil, where $N=2M$, with fundamental stencil $\{a_{-n},a_n\}_{n=1}^M$, consider the resulting $(N+1)$-point odd-centered fundamental stencil given by $\{a_n\}_{n=-M}^M$ where $a_0=0$.
Let $\tilde \FD^k f(\xstar)$ denote the $(N+1)$-point FD approximation
\begin{equation}\label{eq:fake}
	\tilde\FD^k f(\xstar) := \sum_{n=-M}^M c_nf(\xstar+a_nh).
\end{equation}
We first show that there exists a unique set of skew-symmetric coefficients $\{c_n\}_{n=-M}^M$ with $c_0=0$ for which, according to Theorem~\ref{thm:stand_conv}, $\tilde\FD^k f(\xstar)$ is an approximation of order at least $(N+1)-k$.
Since $c_0=0$, this implies that $\tilde \FD^k f\equiv D^kf$.
In other words, the $(N+1)$-point FD approximation in reality has $N$ stencil points, thus implying superconvergence.
Skew-symmetry of the coefficients will come into play when showing the accuracy order does not exceed $N-k+1$.

Assume coefficients $\{c_n\}_{n=-M}^M$ are skew-symmetric, and $c_0=0$ for the ($N+1$)-point, odd-centered FD approximation $\tilde\FD^kf(\xstar)$ in (\ref{eq:fake}).
Similar to Theorem~\ref{thm:stand_conv}, for an $(N+1)$-point stencil, we can eliminate the first $N+1$ terms of the Taylor expansion of the error $E(h)$ by imposing a system of $N+1$ equations on the coefficients $c_n$ of the form
\begin{equation}\label{eq:big}
	A\mathbf c = \mathbf e_{k+1}
\end{equation}
where $\mathbf e_{k+1}\in \mathbb R^{N+1}$ and
\begin{equation}\label{eq:MatAOE}
	A =
	\begin{bmatrix}
		1 &\cdots & 1 &1 & 1 & \cdots & 1\\
		a_{-M}& \cdots& a_{-1}& a_0 & a_1 & \cdots & a_M\\
		a_{-M}^2& \cdots& a_{-1}^2& a_0^2 & a_1^2 & \cdots & a_M^2\\
		\vdots & \ddots & \vdots & \vdots & \vdots & \ddots & \vdots\\
		a_{-M}^{N}&\cdots & a_{-1}^{N}& a_0^N & a_1^{N} & \cdots & a_M^{N}
	\end{bmatrix}\in\mathbb R^{(N+1)\times (N+1)}, \;
	\mathbf c = \begin{bmatrix}
		\tilde c_{-M}\\  \vdots \\ \tilde c_{-1}\\ \tilde c_0\\ \tilde c_1\\ \vdots \\ \tilde c_{M}
	\end{bmatrix}\in\mathbb R^{N+1}.
\end{equation}
Recall, $c_n$ and $\tilde c_n$ are related via (\ref{eq:tilde c}).

Consider the $i^{th}$ component of the product $A\mathbf c$, which we can express as follows given our choice of indexing:
\begin{equation}\label{eq:Ac_i}
    [A\mathbf c]_i 
    = \sum_{n=-M}^M \tilde c_n a_n^{i-1}
    = \tilde c_0 a_0^{i-1} 
    + \sum_{n=1}^M \Big( \tilde c_{-n}a_{-n}^{i-1}+\tilde c_n a_{n}^{i-1}\Big)
\end{equation}
for $i=1,2,...,N+1$.
Then, from the skew-symmetry of the coefficients ($\tilde c_{-n}=-\tilde c_n$), $\tilde c_0 = 0$, and the centeredness of the fundamental stencil ($a_{-n}=-a_n$ and $a_0=0$), (\ref{eq:Ac_i}) implies 
\begin{subequations}
\begin{align}
	[A\mathbf c]_i 
	&= \sum_{n=1}^M \Big( -\tilde c_{n}a_{n}^{i-1}+\tilde c_n a_{n}^{i-1}\Big) = 0, \quad \text{for odd } i, \label{eq:Ac=0}\\
	[A\mathbf c]_i 
	&= 2\sum_{n=1}^M \tilde c_n a_{n}^{i-1}, \quad \text{for even } i. \label{eq:Ac=2}
\end{align}
\end{subequations}

From (\ref{eq:Ac=0}) we can infer that the equations with odd indexes in the original system (\ref{eq:big}) are automatically satisfied.
Using (\ref{eq:Ac=2}) we arrive at a smaller system of equations to determine $\{\tilde c_n\}_{n=1}^M$:
\begin{equation}\label{eq:small}
    \hat A \hat{\mathbf{c}}=\frac{1}{2}\mathbf e_{\ell+1}
\end{equation}
where $\mathbf e_{\ell+1}\in\mathbb R^{M}$ (recall $k=2\ell+1$) and
\[
    \hat A = 
    \begin{bmatrix} 
        a_1 & a_{2} & \cdots & a_{M} \\
        a_1^3 &a_2^3 & \cdots & a_M^3\\
        \vdots & \vdots &\ddots &\vdots & \\
        a_1^{2M-1} &a_2^{2M-1} & \cdots & a_M^{2M-1}\\  
    \end{bmatrix}\in\mathbb R^{M\times M}, \;
    \hat{\mathbf c} =
    \begin{bmatrix}
        \tilde c_1\\
        \tilde c_2\\
        \vdots\\
        \tilde c_M \\
    \end{bmatrix}\in\mathbb R^M.
\]
Note that $\hat A$ is a generalized Vandermonde matrix of the form
\[
	[\hat A]_{ij} = a_j^{z_i},
\]
with $0<a_1<a_2<\cdots < a_M$, where $z_i = 2i-1$ for $i=1,2,...,M$ and thus
\[
	0 \le z_1 < z_2 < \cdots < z_M.
\]
It follows from Lemma~\ref{lem:genVander} that $\hat A$ is invertible.
From the unique solution to reduced system (\ref{eq:small}), we can construct a unique set of skew-symmetric coefficients $\{c_n\}_{n=-M}^M$ so that the original system (\ref{eq:big}) is satisfied.

To show that the accuracy order is exactly $N+1-k$, we look at the expansion of the error $E(h)$ up to $N+2$ terms, after having eliminated the first $N+1$ terms:
\begin{equation*}
    E(h) = \left| K_1h^{N+1-k} + K_2(h)h^{N+2-k} \right|
\end{equation*}
where
\[
	K_1 = \frac{k!}{(N+1)!}f^{(N+1)}(\xstar) \sum_{n=-M}^M \tilde c_n a_n^{N+1}
\]
is independent of $h$, and
\[
	K_2(h) = \frac{k!}{(N+2)!} \sum_{n=-M}^{M} f^{(N+2)}(\xi_n)\tilde c_na_n^{N+2}.
\]
The dependency of $K_2$ on $h$ is implicit via the $\xi_n$.
Using Lemma~\ref{lem:BigTheta2}, it suffices to show that $K_1\neq 0$ and $K_2(h)$ is bounded for small enough $h>0$.
In fact, if $h$ is small enough so that the resulting stencil in contained in $(a,b)$ then we can bound $K_2(h)$ as follows:
\[
	|K_2(h)| \le \frac{k!}{(N+2)!}\sum_{n=-M}^M |\tilde c_na_n^{N+2}| \; \max_{x\in(a,b)}\left|f^{(N+2)}(x)\right|.
\]

Assuming $f^{(N+1)}(\xstar)\neq 0$, using skew-symmetry of coefficients, centeredness of the stencil, and the fact that $N$ is even, we have that $K_1\neq 0$ if and only if
\begin{equation}\label{eq:cr}
    \sum_{n=-M}^M \tilde c_n a_n^{N+1} 
    = 2\sum_{n=1}^M \tilde c_n a_n^{N+1}
    = 2 \hat{\mathbf c}^T \mathbf \rstar
    \neq 0,
\end{equation}
where
\[
	\mathbf \rstar =
	\begin{bmatrix}
	a_1^{N+1} &a_2^{N+1} & \cdots &a_M^{N+1}
	\end{bmatrix}^T.
\] 
Recall that the matrix $\hat A\in\R^{M\times M}$ was shown to be invertible, and hence $\{\mathbf r_n\}_{n=1}^M$ forms a basis for $\R^M$, where $\mathbf r_n$ is the $n^{th}$ row vector of $\hat A$.
It follows that there exists a unique $\mathbf d\in\R^M$ such that
\[
	\mathbf \rstar = \hat A^T\mathbf d
\]
Moreover, since $\hat A\hat{\mathbf c}=\mathbf e_{\ell+1}$, we have
\[
	\hat{\mathbf c}^T\mathbf \rstar = \hat{\mathbf c}^T\hat A^T\mathbf d = \mathbf e_{\ell+1}^T\mathbf d= d_{\ell+1}.
\]
We show that $K_1\neq 0$ by showing that $d_{\ell+1}\neq 0$.

Suppose, as to get a contradiction, that $d_{\ell+1}=0$, which would imply that $\mathbf \rstar$ is in the span of $\{\mathbf r_n\}_{n\neq\ell+1}$.
Construct the matrix $B\in\R^{M\times M}$ by removing the $(\ell+1)^{th}$ row of $\hat A$ and concatenating the row $\mathbf \rstar^T$.
In other words, 
\[
	B = \begin{bmatrix}
		\hat A^{[\ell+1]}\\ \mathbf \rstar^T
	\end{bmatrix}
\]
where $\hat A^{[\ell+1]}\in\R^{(M-1)\times M}$ is the matrix $\hat A$ with its $(\ell+1)^{th}$ row removed.
Note that $B$ is not invertible since its last row is linearly dependent on the previous rows, assuming $d_{\ell+1}=0$.
However, $B$ is also a generalized Vandermonde matrix of the form
\[
	[B]_{ij} = a_j^{z_i}
\]
with $0<a_1<a_2<\cdots < a_M$, where $z_M = N+1$. 
Since
\[
	0<2i-1< N+1, \quad \forall i=1,2,...,M,
\]
this implies $0 \le z_1 < z_2 < \cdots < z_M$.
It follows from Lemma~\ref{lem:genVander} that $B$ must be invertible, and hence our contradiction.
\end{proof}

\end{document}
